% Generated shared reviewed material.
\section{A full-mask Hecke expansion for the type-I piece}
\label{hti:core}
We now continue the signed decomposition through its actual squarefree
inner sum. Fix $K=\mathbb Q(\sqrt{-3})$, the finite-order Hecke character
$\nu$, and the same finite excluded prime set $S$, containing primes above
$6$ and the conductor of $\nu$. For $u\in\mathcal O_K\setminus\{0\}$,
let $\psi_u$ be the primitive Hecke character induced by
$\nu(n)(u/n)_6$ away from $S\cup\operatorname{supp}(u)$.
It is defined before, and independently of, the type-I divisor $d$.
The abelian Kummer extension $K(u^{1/6})/K$ constructs the sextic
factor; multiplying by $\nu$ gives the stated character.
This uses classical abelian reciprocity, not general Artin holomorphy.
Write $\psi_{u,2}$ for the primitive character inducing its square and
$\mathfrak f_u,\mathfrak f_{u,2}$ for the two primitive conductors.
For a squarefree good ideal $d$ put
\[
 T=T(u,d)=S\cup\operatorname{supp}(u)\cup\operatorname{supp}(d),\qquad
 R_T=\prod_{p\in T}Np,\quad Q_u=|D_K|N\mathfrak f_u.
\]
All primes in $T$ are deleted even if absent from either primitive conductor.

\begin{theorem}[Restoration and a uniform scalar contour bound]
\label{hti:theorem}
For $\Re s>1$ the original type-I inner Dirichlet series is
\begin{align}
 F_{u,d}(s)
 &=\sum_{c\ {\rm squarefree}\atop(c,T)=1}
       \frac{\psi_u(c)}{(Nc)^s}
   =\frac{L(s,\psi_u)}{L(2s,\psi_{u,2})}\,C_T(s),\label{hti:ratio}\\
 C_T(s)
 &=\frac{\displaystyle\prod_{p\in T,\ p\nmid\mathfrak f_u}
             (1-\psi_u(p)(Np)^{-s})}
          {\displaystyle\prod_{p\in T,\ p\nmid\mathfrak f_{u,2}}
             (1-\psi_{u,2}(p)(Np)^{-2s})}.
 \label{hti:restore}
\end{align}
It continues meromorphically to $\Re s>1/2$, with only a possible
simple pole at $s=1$ when $\psi_u=1$. In that case its residue is
\[
 \rho_T=\frac{\operatorname*{res}_{s=1}\zeta_K(s)}{\zeta_K(2)}
             \prod_{p\in T}(1+(Np)^{-1})^{-1}.
\]
Fix $0<\delta<1/2$, $\varepsilon>0$ and
$W\in C_c^\infty((0,\infty))$ supported in $[1,2]$. Set
\[\mathcal M_4(W)=\sum_{j=0}^4\sup_x|(x\partial_x)^jW(x)|.\]
Uniformly for $Y\ge1$, $u$ and $d$,
\begin{align}
 \sum_{c\ {\rm squarefree}\atop(c,T)=1}\psi_u(c)W(Nc/Y)
 ={}&\mathbf1_{\psi_u=1}\rho_TY\widehat W(1)\notag\\
 &+O_{\nu,S,\delta,\varepsilon}\!\left(
 \mathcal M_4(W)Y^{1/2+\delta}
 Q_u^{1/4-\delta/2+\varepsilon}R_T^\varepsilon\right).
 \label{hti:scalar}
\end{align}
Here $\widehat W(s)=\int_0^\infty W(x)x^{s-1}\,dx$.
The continuation and functional equation of degree-one Hecke functions
are classical external inputs; the contour argument is given below.
\end{theorem}

At a good prime with $v_p(u)=3$, $\psi_u$ is ramified while its square
is unramified; the second product in~\eqref{hti:restore} must restore a
factor invisible to squaring the zero-extended primitive value.
At $v_p(u)=6$ both characters can forget the prime, but the original row
still vanishes on columns divisible by it. In particular, for principal
$\nu$, $u=\pi^3$ and $u=\pi^6$ at $N\pi=7$, the deleted local factors are
\[
\begin{array}{c|ccc}
 &\text{primitive ratio at }s=1&\text{restoration}&\text{full masked factor}\\\hline
v_p(u)=3&48/49&49/48&1\\
v_p(u)=6&8/7&7/8&1
\end{array}
\]
These are exact local calculations, not numerical evidence about zeros.

For $1\le U$, $U^2\le D$, insert $Y=D/Nd$ into~\eqref{hti:scalar}
and retain the actual $\beta_U(d)$ and outer $\nu(d)\chi_d(u)$.
Classical tame Kummer conductors give $Q_u\ll_\nu Nu$. Thus the normalized
$I_u$ equals its displayed principal-residue sum plus an error $E_u$ with
\begin{equation}\label{hti:outer}
 |E_u|\ll_{\nu,S,\delta,\varepsilon}\mathcal M_4(W)
 D^\delta U^{1-2\delta+\varepsilon}
 (Nu)^{1/4-\delta/2+\varepsilon}.
\end{equation}
The resulting elementary energy envelope is only
\begin{equation}\label{hti:energy}
 \sum_{0<Nu\le H}|E_u|^2
 \ll\mathcal M_4(W)^2D^{2\delta}U^{2-4\delta+\varepsilon}
 H^{3/2-\delta+\varepsilon}.
\end{equation}
It does not prove the short-row target of order $H$ up to small losses.
This is an insufficient upper bound, not a lower bound on the actual
error or a counterexample to coefficient-sensitive averaging. The
original signed mean square and positive determinant bridge stay open.
